\subsection{প্রয়োজন অনুযায়ী Network-এর নকশা}
\label{ch2:c09-10}
সব ক্ষেত্রে একই টপোলজি শ্রেষ্ঠ নয়। নির্ভরযোগ্যতা, বাজেট, দূরত্ব, নোড ও পোর্টের সংখ্যা, ত্রুটি খুঁজে পাওয়ার সুবিধা এবং ভবিষ্যৎ সম্প্রসারণ বিবেচনা করে নকশা বেছে নিতে হয়।
\begin{tabularx}{\linewidth}{|p{.28\linewidth}|p{.19\linewidth}|X|}\hline
\textbf{প্রয়োজন} & \textbf{সম্ভাব্য নকশা} & \textbf{সিদ্ধান্তের কারণ ও সীমা}\\\hline
সাধারণ কলেজের ল্যাব & সুইচভিত্তিক Star & পৃথক প্রান্তিক লিংকে ত্রুটি খুঁজে পাওয়া ও সম্প্রসারণ সহজ; কেন্দ্রের পোর্ট ও বহনক্ষমতা যথেষ্ট হতে হয়।\\\hline
বহু বিভাগের ক্যাম্পাস & Tree/Hybrid & শাখা ভাগ করে পরিচালনা করা যায়; গুরুত্বপূর্ণ কেন্দ্রে বিকল্প পথ ও বিদ্যুতের ব্যবস্থা প্রয়োজন হতে পারে।\\\hline
জরুরি সার্ভার বা ব্যাকবোন & Mesh/Partial Mesh & লিংক নষ্ট হলে বিকল্প পথে যোগাযোগের সুযোগ; বাড়তি কেবল ও পোর্টের ব্যয় বিবেচ্য।\\\hline
পুরোনো ভাগাভাগি করা কেবল & Bus & কেবল কম লাগে; তবে ব্যাকবোন, টার্মিনেশন ও ভাগাভাগি করা ট্রাফিকের সীমা থাকে।\\\hline
পালাক্রমে প্রবেশাধিকার & Token Ring মডেল & টোকেনের অনুমতিতে ডেটা পাঠানো হয়। ত্রুটির প্রভাব নির্দিষ্ট পুনরুদ্ধারব্যবস্থার ওপর নির্ভর করে।\\\hline
\end{tabularx}
\begin{examplebox}{কলেজের Lab ও জরুরি Server একইভাবে যুক্ত হবে কি?}
২০টি ল্যাব PC-এর Star-এ ২০টি প্রান্তিক লিংক ও অন্তত ২০টি উপযুক্ত প্রান্তিক পোর্ট দরকার। ইন্টারনেট বা বহির্মুখী সংযোগের জন্য বাড়তি পোর্ট লাগতে পারে। একই ২০টি PC-এর Full Mesh করলে অংশ \ref{ch2:c09-08}-এর সূত্রে $20\times19/2=190$ লিংক ও ৩৮০টি পোর্টপ্রান্ত লাগে। সাধারণ ল্যাবের জন্য এই ব্যয় অপ্রয়োজনীয় হতে পারে। অল্প কয়েকটি জরুরি সার্ভারে বিকল্প সংযোগের বাড়তি ব্যয় যুক্তিযুক্ত হতে পারে। ক্যাম্পাসের বিভাগগুলো স্তরভিত্তিকভাবে যুক্ত রেখে গুরুত্বপূর্ণ ব্যাকবোনে বিকল্প সংযোগ যোগ করা কার্যকর Hybrid নকশার উদাহরণ।
\end{examplebox}

\subsection{MCQ অনুশীলন}
\begin{multicols}{2}\begin{enumerate}
\item\chTwoMCQBlock{নেটওয়ার্ক টপোলজি কী নির্দেশ করে?
\par ক. সংযোগের বিন্যাস \quad খ. শুধু ইন্টারনেটের গতি \par গ. ফাইলের আকার \quad ঘ. সফটওয়্যারের সংস্করণ }
\item\chTwoMCQBlock{প্রচলিত Ethernet LAN-এর স্টারে কেন্দ্রীয় ডিভাইস সাধারণত কোনটি?
\par ক. সুইচ \quad খ. প্রিন্টার \par গ. স্ক্যানার \quad ঘ. মনিটর }
\item\chTwoMCQBlock{বেশি সরাসরি লিংক ও বিকল্প পথ রাখায় কোন বিন্যাসে কেবল ও পোর্টের খরচ বেশি হতে পারে?
\par ক. Full Mesh \quad খ. Bus \par গ. সাধারণ Ring \quad ঘ. সাধারণ Star }
\end{enumerate}\end{multicols}
\textbf{উত্তরমালা:} ১. ক, ২. ক, ৩. ক।
